\documentclass{article}

\usepackage{amsmath}
\usepackage{amsthm}
\usepackage{amssymb}
\usepackage[margin=1in]{geometry}
\usepackage{hyperref}

\newtheorem{definition}{Definition}
\newtheorem{lemma}{Lemma}
\newtheorem{theorem}{Theorem}

\newcommand{\pref}{\preceq}

\title{Likelihood-Update Merge}
\author{Oscar Laird}
\date{April 2026}

\begin{document}
\maketitle

\section{Overview}


\section{Definitions}

Fix an alphabet $\Sigma$ with a distinguished start symbol $^\wedge$.
Write $\Sigma^\star$ for the finite strings over $\Sigma$, and write
$u \pref x$ when $u$ is a prefix of $x$.

\begin{definition}[Rooted likelihood update]
A \emph{rooted likelihood update} $U$ consists of:
\begin{itemize}
\item a finite set $\mathrm{nodes}(U) \subseteq \Sigma^\star$,
\item a weight $\ell_U(n) \in \mathbb{R}$ for each $n \in \mathrm{nodes}(U)$,
\item the requirement that $^\wedge \in \mathrm{nodes}(U)$,
\item the requirement that $\mathrm{nodes}(U)$ is prefix-closed.
\end{itemize}
\end{definition}

The Rust implementation stores an explicit leaf bit.
Mathematically, this bit is redundant: a present node is a leaf exactly when it
has no present children.

\begin{definition}[Truncation]
Let $U$ be a rooted likelihood update and let $x \in \Sigma^\star$ extend the
start symbol. Define
\[
\tau_U(x) := \max\nolimits_{\pref}\{n \in \mathrm{nodes}(U) : n \pref x\},
\]
the longest present prefix of $x$ in $U$.
This is well defined because $^\wedge \in \mathrm{nodes}(U)$.
\end{definition}

\begin{definition}[Lookup through truncation]
For a rooted likelihood update $U$ and a string $x$, define
\[
L_U(x) := \ell_U(\tau_U(x)).
\]
\end{definition}

\begin{definition}[Semantic merged lookup]
Given rooted likelihood updates $U_1,\dots,U_k$, define their semantic merged
lookup by
\[
L^\ast(x) := \sum_{i=1}^k L_{U_i}(x).
\]
\end{definition}

This function $L^\ast$ is the mathematical specification of the merge.
The formal task is to show that the constructed merged update realizes this
specification.

\section{Construction}

For each string $x \in \Sigma^\star$, define
\[
h(x) := \text{the longest among } \tau_{U_1}(x),\dots,\tau_{U_k}(x).
\]
Since each $\tau_{U_i}(x)$ is a prefix of $x$, this longest string is well
defined.

We define the merged rooted likelihood update $M$ as follows.

\begin{definition}[Merged node set]
A string $n$ belongs to $\mathrm{nodes}(M)$ if and only if there exists a string
$x$ such that
\[
n \pref h(x)
\]
and such that $n \in \mathrm{nodes}(U_i)$ for at least one index $i$.
\end{definition}

\begin{definition}[Merged weights]
For each $n \in \mathrm{nodes}(M)$, define
\[
\ell_M(n) := \sum_{i=1}^k \ell_{U_i}(\tau_{U_i}(n)).
\]
\end{definition}

Thus the weight at a merged node is defined directly from the input updates:
each update contributes the weight of its own truncation point on the branch
leading to $n$.

The merge routine \texttt{merge\_many} is intended to compute exactly this
mathematical object.

\section{Correctness}

\begin{lemma}[Location of the merged truncation point]
Fix $x \in \Sigma^\star$.
For each $i$, let $\tau_i := \tau_{U_i}(x)$.
Then $\tau_M(x)$ is the longest among the strings
$\tau_1,\dots,\tau_k$.
\end{lemma}

\begin{proof}
Each $\tau_i$ is a prefix of $x$, so the strings $\tau_1,\dots,\tau_k$ are
totally ordered by the prefix relation, equivalently by length.

Fix a prefix $n \pref x$.
The update $U_i$ has already stopped by the time the traversal reaches $n$ if
and only if $\tau_i \pref n$.
Therefore all updates have stopped by the time the traversal reaches $n$ if and
only if every $\tau_i \pref n$.

The merged truncation point is the first prefix of $x$ with that property.
Since the $\tau_i$ are linearly ordered, that first such prefix is exactly the
longest among $\tau_1,\dots,\tau_k$.
\end{proof}

\begin{lemma}[Value at the merged truncation point]
Fix $x \in \Sigma^\star$ and let $h := \tau_M(x)$.
Then
\[
\ell_M(h) = \sum_{i=1}^k \ell_{U_i}(\tau_{U_i}(x)).
\]
\end{lemma}

\begin{proof}
Partition the index set $\{1,\dots,k\}$ into
\[
I_{=}:=\{i : \tau_{U_i}(x)=h\}
\quad\text{and}\quad
I_{<}:=\{i : \tau_{U_i}(x)\prec h\}.
\]

If $i \in I_{=}$, then $U_i$ contains the node $h$, so its contribution appears
as the summand $\ell_{U_i}(\tau_{U_i}(x)) = \ell_{U_i}(h)$.

If $i \in I_{<}$, then $U_i$ stopped strictly above $h$.
Its contribution at $h$ is still the summand
$\ell_{U_i}(\tau_{U_i}(x))$, now coming from a strict ancestor of $h$.

These two classes are disjoint and exhaustive, and each index contributes
exactly once.
Summing the direct and carried contributions yields
\[
\ell_M(h)=\sum_{i=1}^k \ell_{U_i}(\tau_{U_i}(x)).
\]
\end{proof}

\begin{theorem}[Constructed merge realizes the specification]
For every string $x \in \Sigma^\star$,
\[
L_M(x) = L^\ast(x).
\]
Equivalently,
\[
L_M(x) = \sum_{i=1}^k L_{U_i}(x).
\]
\end{theorem}

\begin{proof}
By definition of lookup through truncation,
\[
L_M(x) = \ell_M(\tau_M(x)).
\]
Apply the previous lemma with $h = \tau_M(x)$ to obtain
\[
L_M(x) = \sum_{i=1}^k \ell_{U_i}(\tau_{U_i}(x)).
\]
Finally rewrite each summand as $L_{U_i}(x)$.
\end{proof}

\section{Implementation correspondence}

The preceding theorem is the mathematical statement.
The Rust routine \texttt{merge\_many} is intended to compute the update $M$ just
analyzed.
Its running quantities match the two terms in the proof:
\begin{itemize}
\item \texttt{n\_direct\_l} accumulates the direct contributions;
\item \texttt{n\_proper\_truncated\_l} accumulates the carried contributions;
\item \texttt{n\_edge\_hit\_count == len} expresses that every input update has
stopped on the current branch.
\end{itemize}

\section{Next formal step}

This exact theorem is the right first target for Lean.
The later floating-point theorem should be stated separately: first define the
exact real-valued merge semantics, then compare the \texttt{f32} implementation
against that exact target with an explicit rounding-error bound.

\end{document}
